% Source: https://artofproblemsolving.com/community/c3h2475705_combinatronicsesque_problem
% Stable identity: aops-topic-2475705-q01
The answer is $\boxed{31}$.

\textbf{At least $31$ flips may be needed.} Until the light comes on, every observation is ``off'', so a strategy amounts to a fixed sequence of flips. After $k$ flips at most $k+1$ combinations have been shown (the starting one and one new one per flip). If $k\le30$, at most $31$ of the $32$ combinations have been shown, and if the correct combination is one that has not been shown, the light is still off.

\textbf{$31$ flips always suffice.} Number the switches $1$ to $5$. Define flip sequences by $S_1=(1)$ and
\[S_{j+1}=S_j,\ j+1,\ S_j,\]
so $S_2=1,2,1$, $S_3=1,2,1,3,1,2,1$, and $S_j$ has $2^j-1$ flips. We prove by induction on $j$: starting from any position, performing $S_j$ shows all $2^j$ settings of switches $1,\dots,j$ (the other switches stay fixed). For $j=1$ this is clear. If it holds for $j$, then during the first copy of $S_j$ all $2^j$ settings of switches $1,\dots,j$ appear with switch $j+1$ in its initial state; flipping switch $j+1$ shows one new combination; then the second copy of $S_j$ shows all $2^j$ settings of switches $1,\dots,j$ with switch $j+1$ in the other state. That is all $2^{j+1}$ settings, using $2(2^j-1)+1=2^{j+1}-1$ flips.

So $S_5$, which has $31$ flips,
\[1,2,1,3,1,2,1,4,1,2,1,3,1,2,1,5,1,2,1,3,1,2,1,4,1,2,1,3,1,2,1,\]
shows all $32$ combinations from any starting position, and the light comes on at the latest after the $31$st flip.
